
% GA-PRAC-01
% C2 - Escenario 1: Extractora Rio Lindo
% Responsable: Cordova Carreno Mayra Lucila
% Requiere: \usepackage{longtable}
%           \usepackage{array}

\section{C2. Matriz comparativa contextualizada}

\subsection{Escenario 1: Extractora Río Lindo}

La Extractora Río Lindo requiere un sistema embebido
para controlar temperatura, presión y niveles de
tanques. El proyecto presenta requisitos de seguridad
estables e incertidumbre en sensores, interfaz y
reportes. Su criticidad es muy alta debido a los
riesgos para los operadores y la producción.

\small
\setlength{\tabcolsep}{4pt}
\renewcommand{\arraystretch}{1.35}

\begin{longtable}{
|>{\raggedright\arraybackslash}p{0.16\textwidth}
|>{\raggedright\arraybackslash}p{0.25\textwidth}
|>{\raggedright\arraybackslash}p{0.25\textwidth}
|>{\raggedright\arraybackslash}p{0.25\textwidth}|}
\hline
\textbf{Criterio} &
\textbf{Secuencial (cascada)} &
\textbf{Iterativo e incremental} &
\textbf{Ágil} \\
\hline
\endfirsthead

\hline
\textbf{Criterio} &
\textbf{Secuencial (cascada)} &
\textbf{Iterativo e incremental} &
\textbf{Ágil} \\
\hline
\endhead

\textbf{K1. Estabilidad y cambios} &
Favorece la especificación temprana de temperaturas,
presiones y bloqueos, pero dificulta incorporar
cambios posteriores en sensores e interfaz. &
Permite mantener estables los requisitos de seguridad
y refinar por iteraciones la comunicación con sensores
y los reportes. &
Facilita adaptar interfaz y reportes, pero los cambios
frecuentes no aportan la misma ventaja a las reglas
de seguridad ya definidas. \\
\hline

\textbf{K2. Partes interesadas y disponibilidad} &
Permite planificar reuniones con el jefe de planta
dos mañanas por semana y obtener aprobaciones
formales del responsable de seguridad. &
Permite programar revisiones con el jefe de planta
y operadores de distintos turnos al finalizar
cada iteración. &
La colaboración frecuente resulta difícil por la
disponibilidad limitada del jefe de planta y los
turnos rotativos de los operadores. \\
\hline

\textbf{K3. Tipo de sistema y criticidad} &
Favorece la trazabilidad de los requisitos de
protección del sistema embebido, aunque puede
retrasar pruebas de integración. &
Permite comprobar progresivamente alarmas, bloqueos
y comunicación con dispositivos, conservando
requisitos de seguridad verificables. &
Facilita ajustar la interfaz, pero necesita controles
rigurosos para evitar que cambios rápidos comprometan
funciones críticas de protección. \\
\hline

\textbf{K4. Riesgo principal} &
La selección de sensores en el tercer mes podría
invalidar decisiones técnicas anteriores y
generar retrabajo. &
Permite verificar tempranamente la integración
con sensores y válvulas, reduciendo incertidumbre
antes de la puesta en marcha. &
Existe riesgo de priorizar mejoras visibles de
interfaz sobre verificaciones críticas si no se
establecen prioridades y criterios de seguridad
estrictos. \\
\hline

\textbf{K5. Documentación y marco normativo o contractual} &
Favorece la aprobación formal de requisitos de
seguridad y la documentación de pruebas exigidas
por la aseguradora. &
Permite conservar versiones de requisitos y
evidencias de pruebas por incremento, manteniendo
la trazabilidad de funciones de protección. &
Requiere complementar la documentación habitual
de trabajo con requisitos de seguridad aprobados
y evidencias verificables de cada función de
protección. \\
\hline

\end{longtable}

\normalsize
